\section{The roofline model}

The roofline model of Williams, Waterman, and Patterson~\cite{williams2009roofline}
bounds the performance a kernel can attain on a machine by the smaller of two
limits. Let $P_{\max}$ be the peak floating point rate in FLOP/s and $B$ the peak
memory bandwidth in byte/s. For a kernel with arithmetic intensity $I$, measured
in FLOP per byte, the attainable performance is

\begin{equation}
  P(I) = \min\!\left(P_{\max},\; I \cdot B\right).
\end{equation}

The first term is a horizontal line, the compute ceiling: no amount of memory
bandwidth lets a kernel exceed the rate at which the arithmetic units retire
operations. The second term is a line through the origin with slope $B$, the
memory ceiling: a kernel of intensity $I$ cannot run faster than the rate at
which memory can deliver its operands. Plotted on log-log axes, the two combine
into the characteristic roof shape.

The two ceilings meet at the ridge point,

\begin{equation}
  I_{\text{ridge}} = \frac{P_{\max}}{B}.
\end{equation}

A kernel whose intensity is below $I_{\text{ridge}}$ is memory bound: it sits
under the sloped part of the roof, and the way to make it faster is to move less
data or move it more efficiently. A kernel above the ridge is compute bound: it
sits under the flat part, and only more arithmetic throughput helps. The ridge
point is a property of the machine, not the kernel, so the same intensity can be
memory bound on one card and compute bound on another.

Arithmetic intensity is where the honesty of the analysis lives. The theoretical
intensity of a kernel counts the minimum bytes a perfect implementation would
move: read each input once, write each output once. The measured intensity uses
the bytes the kernel actually moved through DRAM, which Nsight Compute reports,
and which include cache misses, register spills, and redundant loads that a real
implementation suffers. The gap between the two is most of the interesting story
in this report, so every point on my rooflines is labeled with which byte count
produced its intensity.
